% GENERATED FILE. Edit canonical lessons, then run npm run generate:books.
% canonical-source-hash: fnv1a64:b531ec0382d7a523
% canonical-lessons: GE-C495-erziehen, GE-C495-erziehung, GE-C495-verwoehnen, GE-C495-gehorchen, GE-C495-pubertaet

\chapter{Kinder erziehen}
\label{ch:ge-kinder-erziehen}
\hlchaptermodality{495}

\begin{chapteropening}
\textbf{By the end of this chapter:} \emph{I can talk about bringing up children.}

Complete the last lesson of chapter 495: I can talk about bringing up children.
\end{chapteropening}

\section[erziehen]{erziehen — to bring up, to raise}

\label{lesson:GE-C495-erziehen}



\begin{quote}
Before the new one: say the German for to reconcile, then the German for reconciliation.
\end{quote}

\subsection*{You'll want to know: erziehen}
\textbf{erziehen} — \textquotedblleft{}to bring up, to raise\textquotedblright{}.

\begin{quote}
Sie hat ihre drei Kinder allein erzogen.
\end{quote}

\begin{quote}
\emph{She brought up her three children on her own.}
\end{quote}

Irregular: \emph{erzog, hat erzogen}.

\begin{cousinweb}
\emph{erziehen} is built on \emph{er-} and \emph{ziehen} \textquotedblleft{}to pull, to draw\textquotedblright{}, inherited from Old High German \emph{ziohan}, from Proto-Germanic \emph{teuhaną} (EtymDB).
\end{cousinweb}

\begin{grammarlens}[title={What you've built}]
One more word for telling, explaining and describing at B1.
\end{grammarlens}

\subsection*{Your turn}
\begin{itemize}
\raggedright
  \item \emph{Say it:} \emph{erziehen}
  \item \emph{Say it:} \emph{erziehen}, once more
  \item \emph{Say it:} a sentence of your own with \emph{erziehen}
  \item \emph{Recall:} say the German for to reconcile, then the German for reconciliation, then say \emph{erziehen} again
\end{itemize}

\begin{tcolorbox}[breakable,colback=teal!4,colframe=teal!35!black,title={Before you move on}]
What does erziehen mean? (\textquotedblleft{}To bring up, to raise\textquotedblright{}.) Say it once more.
\end{tcolorbox}
\section[die Erziehung]{die Erziehung — upbringing}

\label{lesson:GE-C495-erziehung}



\begin{quote}
Before the new one: say the German for reconciliation, then the German for to bring up.
\end{quote}

\subsection*{You'll want to know: die Erziehung}
\textbf{die Erziehung} — \textquotedblleft{}upbringing\textquotedblright{}.

\begin{quote}
Die Eltern streiten oft über die Erziehung.
\end{quote}

\begin{quote}
\emph{The parents often argue about how to bring up the children.}
\end{quote}

No plural. School learning is \emph{die Bildung}.

\begin{cousinweb}
\emph{Erziehung} is built on \emph{erziehen} with \emph{-ung} (EtymDB lists both); \emph{ziehen} goes back to Proto-Germanic \emph{teuhaną} \textquotedblleft{}to pull, to lead\textquotedblright{}.
\end{cousinweb}

\begin{grammarlens}[title={What you've built}]
One more word for telling, explaining and describing at B1.
\end{grammarlens}

\subsection*{Your turn}
\begin{itemize}
\raggedright
  \item \emph{Say it:} \emph{die Erziehung}
  \item \emph{Say it:} \emph{die Erziehung}, once more
  \item \emph{Say it:} a sentence of your own with \emph{die Erziehung}
  \item \emph{Recall:} say the German for reconciliation, then the German for to bring up, then say \emph{die Erziehung} again
\end{itemize}

\begin{tcolorbox}[breakable,colback=teal!4,colframe=teal!35!black,title={Before you move on}]
What does die Erziehung mean? (\textquotedblleft{}Upbringing\textquotedblright{}.) Say it once more.
\end{tcolorbox}
\section[verwöhnen]{verwöhnen — to spoil, to pamper}

\label{lesson:GE-C495-verwoehnen}



\begin{quote}
Before the new one: say the German for to bring up, then the German for upbringing.
\end{quote}

\subsection*{You'll want to know: verwöhnen}
\textbf{verwöhnen} — \textquotedblleft{}to spoil, to pamper\textquotedblright{}.

\begin{quote}
Die Großeltern verwöhnen ihre Enkel mit Süßigkeiten.
\end{quote}

\begin{quote}
\emph{The grandparents spoil their grandchildren with sweets.}
\end{quote}

\emph{Verwöhnt} — spoilt. \emph{Sich verwöhnen lassen} — to let yourself be pampered.

\begin{cousinweb}
\emph{verwöhnen} is usually traced to Middle High German \emph{verwenen} \textquotedblleft{}to accustom wrongly\textquotedblright{}, built on \emph{ver-} and the root of \emph{gewöhnen} \textquotedblleft{}to get used to\textquotedblright{}; EtymDB lists the word without a source.
\end{cousinweb}

\begin{grammarlens}[title={What you've built}]
One more word for telling, explaining and describing at B1.
\end{grammarlens}

\subsection*{Your turn}
\begin{itemize}
\raggedright
  \item \emph{Say it:} \emph{verwöhnen}
  \item \emph{Say it:} \emph{verwöhnen}, once more
  \item \emph{Say it:} a sentence of your own with \emph{verwöhnen}
  \item \emph{Recall:} say the German for to bring up, then the German for upbringing, then say \emph{verwöhnen} again
\end{itemize}

\begin{tcolorbox}[breakable,colback=teal!4,colframe=teal!35!black,title={Before you move on}]
What does verwöhnen mean? (\textquotedblleft{}To spoil, to pamper\textquotedblright{}.) Say it once more.
\end{tcolorbox}
\section[gehorchen]{gehorchen — to obey}

\label{lesson:GE-C495-gehorchen}



\begin{quote}
Before the new one: say the German for upbringing, then the German for to spoil.
\end{quote}

\subsection*{You'll want to know: gehorchen}
\textbf{gehorchen} — \textquotedblleft{}to obey\textquotedblright{}.

\begin{quote}
Der Hund gehorcht nur meiner Mutter.
\end{quote}

\begin{quote}
\emph{The dog only obeys my mother.}
\end{quote}

With a dative: \emph{dem Lehrer gehorchen}. Inseparable: \emph{hat gehorcht}.

\begin{cousinweb}
\emph{gehorchen} is built on \emph{horchen} \textquotedblleft{}to listen, to hearken\textquotedblright{} (EtymDB) with \emph{ge-}: to listen and do as you are told.
\end{cousinweb}

\begin{grammarlens}[title={What you've built}]
One more word for telling, explaining and describing at B1.
\end{grammarlens}

\subsection*{Your turn}
\begin{itemize}
\raggedright
  \item \emph{Say it:} \emph{gehorchen}
  \item \emph{Say it:} \emph{gehorchen}, once more
  \item \emph{Say it:} a sentence of your own with \emph{gehorchen}
  \item \emph{Recall:} say the German for upbringing, then the German for to spoil, then say \emph{gehorchen} again
\end{itemize}

\begin{tcolorbox}[breakable,colback=teal!4,colframe=teal!35!black,title={Before you move on}]
What does gehorchen mean? (\textquotedblleft{}To obey\textquotedblright{}.) Say it once more.
\end{tcolorbox}
\section[die Pubertät]{die Pubertät — puberty}

\label{lesson:GE-C495-pubertaet}



\begin{quote}
Before the new one: say the German for to spoil, then the German for to obey.
\end{quote}

\subsection*{You'll want to know: die Pubertät}
\textbf{die Pubertät} — \textquotedblleft{}puberty\textquotedblright{}.

\begin{quote}
In der Pubertät hatte mein Sohn oft schlechte Laune.
\end{quote}

\begin{quote}
\emph{During puberty my son was often in a bad mood.}
\end{quote}

No plural.

\begin{cousinweb}
\emph{Pubertät} is a learned borrowing from Latin \emph{pūbertās} \textquotedblleft{}the age of maturity\textquotedblright{}, built on \emph{pūbēs} \textquotedblleft{}grown up\textquotedblright{} (Lewis \& Short).
\end{cousinweb}

\begin{grammarlens}[title={What you've built}]
One more word for telling, explaining and describing at B1.
\end{grammarlens}

\subsection*{Your turn}
\begin{itemize}
\raggedright
  \item \emph{Say it:} \emph{die Pubertät}
  \item \emph{Say it:} \emph{die Pubertät}, once more
  \item \emph{Say it:} a sentence of your own with \emph{die Pubertät}
  \item \emph{Recall:} say the German for to spoil, then the German for to obey, then say \emph{die Pubertät} again
\end{itemize}

\begin{tcolorbox}[breakable,colback=teal!4,colframe=teal!35!black,title={Before you move on}]
What does die Pubertät mean? (\textquotedblleft{}Puberty\textquotedblright{}.) Say it once more.
\end{tcolorbox}
